\documentclass[10pt]{article}
\usepackage[T1]{fontenc}
\usepackage[utf8]{inputenc}
\usepackage[french]{babel}
\usepackage{amsmath,amssymb,mathrsfs,mathpazo}
\usepackage{geometry}
\geometry{a4paper,margin=2cm}
\usepackage{graphicx,array,tabularx,tabularray,colortbl,multirow,multicol,booktabs,enumitem,pifont,ulem,fancybox,diagbox,cancel}
\usepackage{pstricks,pstricks-add,pst-plot,pst-tree,pst-node,pst-text,pst-func,pst-math,pst-bspline,pst-3dplot,pst-eucl,pst-solides3d}
\let\pscalculate\undefined
\usepackage{pst-bezier}
\usepackage{tikz,pgfplots,tkz-tab}
\usetikzlibrary{arrows,arrows.meta,patterns,calc,positioning,trees,decorations.markings,shapes.geometric,shapes.misc,shapes.arrows,intersections,decorations.pathmorphing,decorations.pathreplacing,backgrounds,shadows,babel}
\pgfplotsset{compat=1.18}
\usepgfplotslibrary{fillbetween}
\usepackage[np]{numprint}
\usepackage{listings,makecell,pas-tableur,twoopt}
\let\barre\undefined
\usepackage[pstricks]{tablvar}
\definecolor{darkgreen}{RGB}{0,100,0}
\providecommand{\red}{\color{red}}
\providecommand{\blue}{\color{blue}}
\providecommand{\green}{\color{green}}
\providecommand{\black}{\color{black}}
\providecommand{\white}{\color{white}}
\providecommand{\gray}{\color{gray}}
\providecommand{\cyan}{\color{cyan}}
\providecommand{\magenta}{\color{magenta}}
\providecommand{\yellow}{\color{yellow}}
\providecommand{\violet}{\color{violet}}
\providecommand{\orange}{\color{orange}}
\providecommand{\euro}{\texteuro}
\providecommand{\EUR}{\texteuro}
\providecommand{\ieme}{\textsuperscript{e}}
\providecommand{\degres}{\ensuremath{^\circ}}
\providecommand{\degree}{\ensuremath{^\circ}}
\providecommand{\vv}[1]{\overrightarrow{#1}}
\providecommand{\Coord}[1]{#1}
\providecommand{\decosix}{}
\providecommand{\starredbullet}{\ensuremath{\bullet}}
\pagestyle{empty}
\setlength{\parindent}{0pt}
\usepackage{adjustbox}\DeclareUnicodeCharacter{2236}{:}\DeclareUnicodeCharacter{2190}{\ensuremath{\leftarrow}}
\providecommand{\R}{\mathbb{R}}
\providecommand{\N}{\mathbb{N}}
\providecommand{\Z}{\mathbb{Z}}
\providecommand{\Q}{\mathbb{Q}}
\providecommand{\D}{\mathbb{D}}
\providecommand{\C}{\mathbb{C}}
\providecommand{\e}{{\mathrm{e}}}
\providecommand{\vect}[1]{\overrightarrow{#1}}
\providecommand{\vectt}[1]{\overrightarrow{\mathrm{#1}}}
\providecommand{\vvt}[1]{\overrightarrow{\mathrm{#1}}}
\providecommand{\np}[2][]{{#2~#1}}
\providecommand{\vv}[1]{{\overrightarrow{#1}}}
\renewcommand{\d}{{\mathrm{d}}}
\providecommand{\degres}{{^\circ}}
\providecommand{\degre}{{^\circ}}
\providecommand{\blue}{\color{blue}}
\providecommand{\red}{\color{red}}
\providecommand{\enskip}{\quad}
\providecommand{\llbracket}{\lbrack\!\lbrack}
\providecommand{\rrbracket}{\rbrack\!\rbrack}
\providecommand{\euro}{{\text{€}}}
\providecommand{\no}{{\text{n°}}}
\providecommand{\interval}[2]{{[#1;#2]}}
\providecommand{\intervFO}[2]{{[#1;#2[}}
\providecommand{\dotfill}{\ldots}

\providecommand{\ds}{\displaystyle}
\providecommand{\textding}[1]{#1}
\providecommand{\Oij}{\(\mathrm{O};\vec{\imath},\vec{\jmath}\)}
\providecommand{\Oijk}{\(\mathrm{O};\vec{\imath},\vec{\jmath},\vec{k}\)}
\providecommand{\Ouv}{\(\mathrm{O};\vec{u},\vec{v}\)}
\providecommand{\intervFF}[2]{\ensuremath{[#1;#2]}}
\renewcommand{\euro}{\ifmmode\text{\texteuro}\else\texteuro\fi}
\DeclareUnicodeCharacter{00B0}{\ifmmode^{\circ}\else\textdegree\fi}
\DeclareUnicodeCharacter{20AC}{\ifmmode\text{\texteuro}\else\texteuro\fi}
% Énoncés : Éducation nationale. Transcription et corrigés : APMEP, Denis Vergès et contributeurs. Découpage et classement : CoopMaths, Sébastien Lozano, Éric Elter et contributeurs. Adaptation de la présentation : Maths Explorer.
% CC BY-SA : https://coopmaths.fr/www/about/
% https://coopmaths.fr/alea/static/eam/2026/tex/eam_specialite_2026_06_ameriquenord_3_cor.tex
\begin{document}
\noindent\textit{APMEP / CoopMaths. CC BY-SA. Présentation adaptée par Maths Explorer.}\par\medskip
\begin{enumerate}
		\item Pour tout $x\in\R$,\quad $f(x)=(4x-4)\e^{-0,5x}+5$.

		On a $f$ de la forme : \quad $f = u \times v + 5$, \quad et donc : $f'=u'v + uv' + 0$.

		On dérive le produit $u\times v$ avec :\quad $u : x\longmapsto 4x-4$, \quad donc \quad $u': x \longmapsto 4$,

		et \quad $v : x\longmapsto \e^{-0,5x}$,\quad donc \quad $v':x\longmapsto -0,5\,\e^{-0,5x}$.

		$\forall x \in \R,\quad \begin{aligned}[t]
			f'(x)&=u'(x)v(x)+u(x)v'(x) + 0\\
			&=4\e^{-0,5x}+(4x-4)\times(-0,5)\e^{-0,5x}\\
			&=\big[4-0,5\times(4x-4)\big]\e^{-0,5x}\\
			&=\big[4-2x+2\big]\e^{-0,5x}\\
			&=(-2x+6)\e^{-0,5x}.
		\end{aligned}$

		On arrive bien à l'expression donnée pour la fonction $f'$.

		\item La fonction exponentielle est à valeurs strictement positives, donc pour tout $x\in\R$,\quad $\e^{-0,5x}>0$.

		Le signe de $f'(x)$ est celui de $-2x+6$.
		$-2x+6=0\Longleftrightarrow x=3$,\quad et \quad $-2x+6>0\Longleftrightarrow x<3$.

		On en déduit donc :
		\begin{itemize}[left=5mm,label={\textbullet~} ]
			\item sur $]-\infty ; 3]$, $f'$ est à valeurs positives, et donc $f$ est croissante;

			\item sur $[3\,;\,+\infty[$, $f'$ est à valeurs négatives et donc $f$ est décroissante.
		\end{itemize}

		On peut donc donner le tableau de variations :

		\hfill~\begin{tikzpicture}
			\tkzTabInit[lgt=3]{$x$/0.7, signe de $f'(x)$/0.7, variations de $f$/2.5}{$-\infty$,$3$,$+\infty$}
			\tkzTabLine{,+,z,-,}
			\tkzTabVar{-/,+/$8\e^{-1,5}+5$,-/}
		\end{tikzpicture}\hfill~

		\item La tangente à $\mathcal{C}_{f}$ est horizontale quand $f'(x)=0$, c'est-à-dire en $x=3$ uniquement.

		$f(3)=(4\times 3-4)\e^{-0,5\times 3}+5=8\e^{-1,5}+5$.

		$\mathcal{C}_{f}$ admet donc un seul point à tangente horizontale.

		Les coordonnées exactes de ce point sont :\quad $\left(3\,;\,8\e^{-1,5}+5\right)$.
	\end{enumerate}

\end{document}
